\chapter{Comparison and generalized rules}
\label{chap:comparison-and-rules}

This chapter begins the internal-proof layer.  It consumes the explicit literature and computation inputs above to prove comparison, page-extension, Leibniz, Moss-adapter, and Mahowald results inside the project.

\section{Extension-SS consumers}

% SOURCE: AIM-7 concrete tracer bullet; aimpaper/main.tex:430-453.
\begin{proposition}[Concrete eta detection slice]
  \label{prop:classical-eta-extension-iff-detected}
  \uses{def:classical-eta-ess-concrete,prop:eta-ess-regression}
  \lean{KIP126.Classical.ExtensionSS.extension_iff_detected}
  \lean{KIP126.Classical.ExtensionSS.FExtension}
  \lean{KIP126.Classical.ExtensionSS.DetectedBy}
  \lean{KIP126.Classical.ExtensionSS.Essential}
  \lean{KIP126.Classical.ExtensionSS.Crossing}
  \notready
  For the concrete finite eta rows, extension is equivalent to membership in
  the set-valued detection predicate.  Essentiality and crossing are defined
  only on this construction, so detection representatives are not forced to
  be unique.
\end{proposition}

\begin{proof}
  Expand the two predicates as membership in the detected and differential
  sets and rewrite with the concrete input's equality between those sets.
  The eta rows themselves are supplied by the catalogued regression evidence.
\end{proof}

% SOURCE: AIM-7 concrete tracer bullet; aimpaper/main.tex:430-453.
\begin{definition}[Concrete classical $\eta$-ESS tracer bullet]
  \label{def:classical-eta-ess-concrete}
  \uses{def:classical-adams-ss,def:ess-complex-shape,thm:external-adams-h4-d2-slice,prop:eta-ess-regression,def:catalogued-external-evidence}
  \lean{KIP126.Classical.ExtensionSS.etaESS}
  \lean{KIP126.Classical.ExtensionSS.E₀}
  \lean{KIP126.Classical.ExtensionSS.abutment}
  \lean{KIP126.Classical.ExtensionSS.ClassicalEtaESSAdapter}
  \notready
  On the AIM-5 concrete classical Adams data, the $\eta$ tracer bullet uses
  Mathlib's `CategoryTheory.SpectralSequence` with page-zero object
  $E_\infty(X)\oplus E_\infty(Y)$, differential degree $(n,n)$, and
  componentwise abutment $\ker(\pi_*\eta)\oplus\operatorname{coker}(\pi_*\eta)$.
  The finite differential rows carry the AIM paper locator and are injected
  through the `etaEssRegression` claim ledger; this does not add a global
  differential axiom.
\end{definition}

% SOURCE: aimpaper/main.tex:327-342, unlabelled Corollary 2.6.
\begin{corollary}[Map filtration forces initial ESS differentials to vanish]
  \label{cor:fess-below-map-filtration-zero}
  \uses{thm:external-map-filtration-factorization,prop:fextension-iff-detected}
  \notready
  If $\operatorname{AF}(f)=k$, then $d_i^f=0$ for every $i<k$.
\end{corollary}

\begin{proof}
  Factor $f$ into $k$ homology-trivial maps.  Each factor raises Adams
  filtration by at least one, so the composite raises it by at least $k$.
  Proposition~\ref{prop:fextension-iff-detected} then excludes a nonzero
  target in any smaller filtration shift.
\end{proof}

\section{Moss convergence adapter}

% SOURCE: aimpaper/main.tex:2387-2477; external Moss convergence theorem is required.
\begin{theorem}[Moss convergence adapter]
  \label{thm:moss-convergence-adapter}
  \uses{def:toda-bracket,def:massey-product,thm:toda-product-identities,def:adams-detection,def:external-result}
  \notready
  Given an explicit external result supplying the relevant Moss convergence
  theorem and hypotheses excluding crossings and indeterminacy, a Massey
  product in the Adams $E_r$-page detecting a permanent class yields a
  corresponding element of the stated Toda bracket in stable homotopy.
\end{theorem}

\begin{proof}
  Unpack the external theorem at the chosen spectra and representatives,
  discharge its convergence-and-detection witness and no-crossing hypotheses
  from the supplied spectral-sequence data, and transport the detected class through the
  associated-graded isomorphism.  Keep the resulting Toda indeterminacy
  explicit rather than choosing a canonical element.
\end{proof}

\section{Synthetic comparison and normalized lifts}

% SOURCE: aimpaper/main.tex:685-691; AIM-6 concrete synthetic comparison slice.
\begin{proposition}[The $h_4$ classical--synthetic $d_2$ comparison]
  \label{prop:aim6-h4-comparison}
  \uses{def:classical-adams-ss,def:synthetic-adams-ss,def:synthetic-tridegree,def:reindexed-ss-morphism,def:external-result}
  \lean{KIP126.Synthetic.SpectralSequence.SyntheticAdamsSS.h₄}
  \lean{KIP126.Synthetic.SpectralSequence.SyntheticAdamsSS.h₀h₃Squared}
  \lean{KIP126.Synthetic.SpectralSequence.SyntheticLambdaAction}
  \lean{KIP126.Synthetic.SpectralSequence.SyntheticAdamsSS.lambdaMap}
  \lean{KIP126.Synthetic.SpectralSequence.SyntheticAdamsSS.lambdaMap_comm}
  \lean{KIP126.Synthetic.SpectralSequence.SyntheticAdamsSS.weightPreserving}
  \lean{KIP126.Comparison.ClassicalSynthetic.sourcePageMap}
  \lean{KIP126.Comparison.ClassicalSynthetic.sourcePageMap_differential_naturality}
  \notready
  The upstream data layer now attaches multiplication by $\lambda$ to each
  Mathlib page through \texttt{SyntheticLambdaAction}; its page maps are
  morphisms of homological complexes, so \texttt{lambdaMap\_comm} follows from
  chain-map commutativity.  The concrete \texttt{fixedWeightPage} and
  \texttt{ReindexedSpectralSequenceMap.differential\_comm} similarly derive
  differential naturality for fixed-weight comparison maps.  The final $h_4$
  correspondence remains `notready`: the project still lacks
  representative-level comparison across the two weights and compatibility
  of that comparison with the typed action.
\end{proposition}

\begin{proof}
  The AIM-5 catalogue wrapper supplies the classical statement and degree
  equalities.  The remaining proof obligation is the representative-induced
  comparison across the two fixed weights and its compatibility with the
  typed action; no proof of the differential correspondence is claimed here.
\end{proof}

% PROOF-SOURCE: BHS Theorem A.1 and its special-fiber edge isomorphism;
% aimpaper/main.tex:813-823.  This is the exponent-one boundary case that is
% not covered by aimpaper/main.tex:849-855.
\begin{proposition}[$E_\infty$ at the first $\lambda$-quotient]
  \label{prop:synthetic-einfty-first-quotient}
  \uses{thm:external-lambda-bockstein,thm:external-synthetic-rigidity,def:synthetic-lambda-quotient,def:synthetic-adams-ss,def:ss-convergence-detection-witness,def:bhs-aim-regrading}
  \notready
  For every admissible $X$, the special-fiber edge map identifies the
  synthetic Adams abutment of $\nu X/\lambda$ with the classical initial page:
  \[
    E_\infty^{s,t,w}(\nu X/\lambda)\cong
    \begin{cases}
      E_2^{s,t}(X),&w=t,\\
      0,&w\ne t.
    \end{cases}
  \]
  Under this isomorphism the edge class has a unique detected homotopy class
  in its tridegree.  This proposition is the explicit exponent-one interface;
  it does not extend the finite $E_\infty$ formula $Z_i/B_j$ of
  Theorem~\ref{thm:external-synthetic-einfty-quotient} outside its stated
  range $r\ge2$.
\end{proposition}

\begin{proof}
  Apply the synthetic Adams $E_2$ functor to the cofiber of $\lambda$.
  Rigidity identifies its page with
  $E_2^{*,*}(X)\otimes\mathbb F_2[\lambda]/(\lambda)$, concentrated in weight
  $w=t$.  The $\lambda$-Bockstein comparison gives the edge isomorphism
  $E_2^{s,t}(X)\cong\pi_{t-s,t}(\nu X/\lambda)$.  Strong convergence and the
  fact that this edge map already exhausts the abutment force every
  differential and hidden additive extension at this quotient to vanish.
  Regrade to AIM's three coordinates to obtain the displayed presentation and
  uniqueness statement.
\end{proof}

% SOURCE: aimpaper/main.tex:865-875, load-bearing compatibility remark after Propositions 3.11-3.12.
\begin{proposition}[Compatibility of $E_\infty$ formulas with $\lambda$ and $\rho$]
  \label{prop:synthetic-einfty-map-compatibility}
  \uses{thm:external-synthetic-einfty-nu,thm:external-synthetic-einfty-quotient,prop:synthetic-einfty-first-quotient,def:synthetic-lambda-quotient,def:bhs-aim-regrading}
  \notready
  For $r'\ge r\ge2$, multiplication is the degree-zero synthetic map
  \[
    \lambda^{r'-r}:
    \Sigma^{0,r-r'}(\nu X/\lambda^r)\longrightarrow\nu X/\lambda^{r'}.
  \]
  Under the preceding isomorphisms, its map on $E_\infty$ is
  \[
    E_\infty^{s,t,w}(\nu X/\lambda^r)
      \longrightarrow
    E_\infty^{s,t,w-r'+r}(\nu X/\lambda^{r'}),
  \]
  and, when the displayed cycle and boundary indices are positive, it is the
  quotient map
  \[
    \frac{Z_{r-t+w}^{s,t}(X)}{B_{1+t-w}^{s,t}(X)}
      \longrightarrow
    \frac{Z_{r-t+w}^{s,t}(X)}{B_{1+t-w+r'-r}^{s,t}(X)}.
  \]
  For $r\ge r''\ge2$, reduction
  $\rho:\nu X/\lambda^r\to\nu X/\lambda^{r''}$
  is, in the same range, the inclusion
  \[
    \frac{Z_{r-t+w}^{s,t}(X)}{B_{1+t-w}^{s,t}(X)}
      \hookrightarrow
    \frac{Z_{r''-t+w}^{s,t}(X)}{B_{1+t-w}^{s,t}(X)}.
  \]
  The project's zero-index convention supplies the remaining cycle- and
  boundary-subscript cases within these exponent ranges.

  For the boundary case $r=1<r'$, the same suspended multiplication map sends
  the special-fiber edge group to the indicated weight of the $r'$ quotient;
  it is the canonical quotient and hence is surjective or zero in each
  tridegree.  For $r''=1<r$, reduction identifies the surviving-cycle subgroup
  with a subgroup of the special-fiber edge group and hence is injective or
  zero.  When both exponents are one the map is the identity.  These boundary
  statements use Proposition~\ref{prop:synthetic-einfty-first-quotient}, not
  the finite-quotient formula outside its range.

  There are also two explicit infinite-endpoint clauses.  For every $n>0$,
  multiplication
  \[
    \lambda^n:\Sigma^{0,-n}\nu X\longrightarrow\nu X
  \]
  is described as follows.  On a tridegree $(s,t,w)$ of the suspended domain,
  put $j=1+t-(w+n)$.  When the source is nonzero, suspension identifies it
  with $Z_\infty^{s,t}(X)/B_j^{s,t}(X)$ and the map is the
  boundary-cutoff quotient
  \[
    Z_\infty^{s,t}(X)/B_j^{s,t}(X)
      \longrightarrow
    Z_\infty^{s,t}(X)/B_{j+n}^{s,t}(X),
  \]
  since the target cutoff is $1+t-w=j+n$.  It is therefore surjective or zero
  in each tridegree.  For reduction
  \[
    \rho:\nu X\longrightarrow\nu X/\lambda^n
  \]
  fix a tridegree $(s,t,w)$ and put $j=1+t-w$.  It is the cycle-cutoff inclusion
  \[
    Z_\infty^{s,t}(X)/B_j^{s,t}(X)
      \hookrightarrow
    Z_{n-t+w}^{s,t}(X)/B_j^{s,t}(X)
  \]
  whenever $n\ge2$ and $0\le t-w<n$, and is injective or zero otherwise.  At
  $n=1$ it is the inclusion of permanent representatives
  into the special-fiber edge group (or zero off its unique weight).  The
  infinite source uses Theorem~\ref{thm:external-synthetic-einfty-nu}; it is
  not obtained by substituting $r=\infty$ into the finite formula.
\end{proposition}

\begin{proof}
  Express both synthetic maps on the $\lambda$-Bockstein exact couple.  The
  first changes the boundary cutoff by $r'-r$ and shifts the weight by
  $r-r'$, while the second changes the cycle cutoff without shifting weight.
  For exponents at least two, transport these maps through the two external
  $E_\infty$ isomorphisms and check the zero-index edge cases separately.  If
  an endpoint exponent is one, use the special-fiber edge presentation: the
  multiplication map quotients the initial-page group by the later boundary
  cutoff, while reduction includes the surviving-cycle subgroup.  This gives
  the asserted surjective/zero and injective/zero alternatives without
  applying the finite-quotient theorem at exponent one.

  For an infinite endpoint, compare the untruncated formula
  $Z_\infty/B_j$ directly with its weight-shifted copy and with the finite
  quotient formula.  Multiplication by $\lambda^n$ raises the boundary cutoff
  from $j$ to $j+n$, hence is the displayed quotient.  Reduction retains the
  boundary cutoff and includes $Z_\infty$ into the finite cycle subgroup; for
  exponent one use the special-fiber edge group.  These calculations prove
  the two infinite clauses independently of any limiting shorthand.
\end{proof}

% SOURCE: aimpaper/main.tex:878-910, Example exam:synEinfty.
\begin{proposition}[Synthetic $14$-stem regression]
  \label{prop:synthetic-14-stem-regression}
  \uses{thm:external-synthetic-rigidity,thm:external-synthetic-einfty-nu,thm:external-synthetic-einfty-quotient,def:external-evidence}
  \notready
  Given the explicit classical evidence
  $d_2(h_4)=h_0h_3^2$,
  $d_3(h_0h_4)=h_0d_0$, and
  $d_3(h_0^2h_4)=h_0^2d_0$, the synthetic spectral sequences of
  $S^{0,0}$, $S^{0,0}/\lambda^2$, and $S^{0,0}/\lambda^3$ have exactly the
  differentials, permanent multiples, and $\lambda$-torsion classes listed in
  Example 3.15.
\end{proposition}

\begin{proof}
  Apply rigidity to insert the factors $\lambda$ and $\lambda^2$, truncate at
  the quotient exponent, and evaluate the two $E_\infty$ formulas.  A finite
  degree check verifies the list and serves as a regression test for grading
  and zero-index conventions.
\end{proof}

% SOURCE: aimpaper/main.tex:946-961, cofiber assertion in Notation not:fhat.
\begin{proposition}[Cofiber of a normalized synthetic map]
  \label{prop:synthetic-hat-cofiber}
  \uses{def:synthetic-hat-map,def:normalized-synthetic-exactness-case,thm:external-nu-cofiber-criterion,thm:external-synthetic-triangle-lift}
  \notready
  Let $X\xrightarrow fY\xrightarrow gC(f)$ be a cofiber sequence equipped with
  a normalized cofiber exactness case.  Let $\widehat f$ and $\widehat g$ be
  the triangle-compatible components obtained simultaneously from the same
  lifted rotation carried by that case.  Then
  \[
    C(\widehat f)\simeq\Sigma^{0,-e(g)}\nu C(f).
  \]
  Thus the tagged zero and monomorphism cases, both of which explicitly have
  $e(g)=0$, give no shift.  The tagged epimorphism case explicitly has
  $e(g)=1$ and gives a $(0,-1)$ shift.  The homology-map adjective alone does
  not select a shift in a degenerate case.  No such equivalence is asserted for
  $\lambda^{k-1}\widetilde f$ formed from an arbitrary full lift.
\end{proposition}

\begin{proof}
  Eliminate the tagged exactness datum.  In the zero constructor, rotate so
  that $f$ is the positive-filtration connecting map; in the monomorphism
  constructor, use the original short exact homology sequence; in the
  epimorphism constructor, rotate so that $g$ is the positive-filtration
  connecting map.
  Apply the cofiber criterion to the other two components in that chosen
  rotation.  Because $\widehat f$ and $\widehat g$ are components of this one
  lifted triangle, its cofiber identifies with the displayed object.  After
  rotating back, use the constructor's explicit $e(g)$ field to obtain exactly
  the asserted shift.  The argument never
  replaces $\widehat f$ by a torsion-equivalent arbitrary full lift.
\end{proof}

% SOURCE: aimpaper/main.tex:965-971, unlabelled Proposition 3.20.
\begin{proposition}[Normalized synthetic triangle]
  \label{prop:normalized-synthetic-triangle}
  \uses{def:synthetic-hat-map,def:normalized-synthetic-exactness-case,prop:positive-adams-filtration-homology-zero,prop:synthetic-hat-cofiber,thm:external-synthetic-triangle-lift,def:cofiber-triangle,def:hf2-homology}
  \notready
  If $X\xrightarrow fY\xrightarrow gZ\xrightarrow h\Sigma X$ is a
  distinguished triangle and $e(f)+e(g)+e(h)=1$, there is a simultaneous
  choice of triangle-compatible normalized lifts that forms a distinguished
  synthetic triangle
  \[
    \nu X\xrightarrow{\widehat f}\Sigma^{0,-e(f)}\nu Y
    \xrightarrow{\widehat g}\Sigma^{0,-e(f)-e(g)}\nu Z
    \xrightarrow{\widehat h}\Sigma^{0,-1}\nu\Sigma X.
  \]
\end{proposition}

\begin{proof}
  The sum-one hypothesis selects exactly one of the three explicit
  $e$-patterns.  Proposition~\ref{prop:positive-adams-filtration-homology-zero}
  makes the unique positive-filtration map zero on homology.  Exactness of the
  cofiber long exact sequence then says respectively that $H_*f$ is zero,
  injective, or surjective, and supplies the corresponding rotated short exact
  sequence.  Transport along the chosen equivalence with the selected cofiber
  triangle, rotate so that this positive-filtration map is the connecting map,
  and apply the triangle-lift theorem once.  These fields form the appropriate
  tagged exactness datum.  Rotate that same synthetic triangle back.  Its three components define the normalized lifts
  simultaneously; Proposition~\ref{prop:synthetic-hat-cofiber} identifies the
  required cofibers.  The equality of the three $e$-values makes the total
  weight shift exactly $-1$.
\end{proof}

\section{Classical--synthetic extension comparison}

% PROOF-SOURCE: the standard obstruction calculation for a countable inverse
% system of abelian torsors, combined with BHS Proposition A.13.
\begin{lemma}[Mittag--Leffler criterion for an inverse-limit ESS extension]
  \label{lem:synthetic-ess-coherent-tower-limit}
  \uses{def:synthetic-ess-solution-tower,prop:lambda-quotient-restriction-coherence,thm:external-synthetic-lambda-complete}
  \notready
  In the setting of
  Definition~\ref{def:synthetic-ess-solution-tower}, suppose every finite
  solution fiber is nonempty and that the compatible maps $F_m$ are the
  reductions of a map $F_\infty$ between $\lambda$-complete synthetic spectra.
  There is a canonical torsor obstruction
  \[
    \omega(\operatorname{Sol}_\bullet)
      \in \varprojlim^{1}_{m} K_m
  \]
  whose vanishing is equivalent to the existence of a coherent solution
  tower.  Either of the following is sufficient for vanishing:
  \begin{enumerate}
    \item every finite solution lifts, that is, each
      $\operatorname{res}_{m+1,m}$ is surjective;
    \item the difference-group tower $(K_m)$ is Mittag--Leffler: for every
      $m$, the images of $K_l\to K_m$ stabilize as $l$ increases.  Surjective
      transition maps are a sufficient special case.  Then
      $\varprojlim^1_m K_m=0$.
  \end{enumerate}
  A coherent solution tower determines the untruncated relation
  \[
    d_n^{F_\infty}(x_\infty)=y_\infty
  \]
  in the complete target coset.  Conversely, every untruncated representative
  solution restricts to a coherent finite tower.
\end{lemma}

\begin{proof}
  Choose one point in each finite solution torsor.  The differences between a
  chosen point and the restriction of the next one form the usual cocycle for
  the inverse system $(K_m)$; changing the choices changes it by a coboundary.
  Its class is $\omega$, and a coboundary correction makes the chosen points
  coherent exactly when $\omega=0$.  Surjective lifting constructs such a
  tower recursively.  The Mittag--Leffler hypothesis gives
  $\varprojlim^1K_m=0$ and hence the same conclusion.

  At the representative-space level, the solution equation is a homotopy fiber
  of the filtered two-term-complex map, and homotopy limits preserve that
  fiber.  The $\lambda$-adic equivalences identify the limit source, target,
  map, and shorter-image corrections with their untruncated counterparts.
  Thus a coherent finite solution is exactly an $F_\infty$ solution; restriction
  proves the converse.
\end{proof}

% SOURCE: aimpaper/main.tex:1025-1048, Proposition prop:9770ae6e.
\begin{proposition}[$\lambda^n$- and $\rho$-ESS have only $d_0$]
  \label{prop:lambda-rho-ess-only-d0}
  \uses{def:lambda-rho-delta-triangle,prop:synthetic-einfty-map-compatibility,prop:synthetic-einfty-first-quotient,prop:synthetic-ess-weightwise-transport,prop:fess-exact-middle-boundary}
  \notready
  In the extension spectral sequences of $\lambda^n$ and $\rho_{n,m}$, the
  only possibly nonzero differentials are
  $d_0^{\lambda^n}=\lambda^n$ and $d_0^\rho=\rho$.  These $d_0$ extensions have
  no crossings.  On each tridegree, $\lambda^n$ is surjective or zero and
  $\rho$ is injective or zero on $E_\infty$.
\end{proposition}

\begin{proof}
  First suppose $m<\infty$.  If one of the quotient exponents is one, use the
  mod-$\lambda$ collapse and unique-detection statement of
  Proposition~\ref{prop:synthetic-einfty-first-quotient}; when both relevant
  exponents are at least two, use the two finite $E_\infty$ formulas.  If a
  tridegree is outside the quotient range its group is zero, while in range
  $\rho$ is a cycle inclusion and $\lambda^n$ is a boundary quotient.

  If $m=\infty$, instead use the two infinite-endpoint clauses of
  Proposition~\ref{prop:synthetic-einfty-map-compatibility}: on the
  untruncated $E_\infty$ formula, $\lambda^n$ is the quotient that raises the
  boundary cutoff by $n$, and
  $\rho:\nu X\to\nu X/\lambda^n$ includes permanent representatives into the
  finite cycle cutoff (using the special-fiber clause when $n=1$).  Thus the
  sequence is exact at its middle term in every tridegree in both the finite
  and infinite branches.
  Apply the synthetic exact-middle theorem to eliminate higher differentials,
  and use the degree-zero filtration shift to obtain no-crossing.
\end{proof}

% SOURCE: aimpaper/main.tex:1068-1131, proof of Proposition prop:6de7d130.
\begin{lemma}[Shorter $\delta_n$ images are classical boundaries]
  \label{lem:delta-images-equal-classical-boundaries}
  \uses{thm:external-lambda-bockstein,prop:synthetic-einfty-map-compatibility,prop:synthetic-ess-weightwise-transport,cor:fess-triangle-shift}
  \notready
  For the infinite connecting map
  $\delta_n:\nu X/\lambda^n\to\Sigma^{1,-n}\nu X$, the subgroup generated by
  extension differentials of length less than $r$ in the target tridegree is
  exactly the image of the classical boundary subgroup
  $B_{r-1}^{s+r,t+r-1}(X)$ under the corresponding $\lambda$-multiple.
\end{lemma}

\begin{proof}
  Induct on $r$.  Write each generator of $B_{r-1}$ as the target of an
  essential classical $d_{r'}$ with $2\le r'<r$.  The Bockstein comparison and
  triangle-shift propagation turn it into a shorter $\delta_n$ extension with
  the required $\lambda$ power.  Conversely, apply the same comparison to
  every shorter $\delta_n$ extension.  The inductive hypothesis removes its
  lower-page indeterminacy.
\end{proof}

% SOURCE: aimpaper/main.tex:1068-1131, Proposition prop:6de7d130, m=infinity case.
\begin{proposition}[Infinite $\delta_n$-ESS formula]
  \label{prop:delta-ess-formula-infinite}
  \uses{lem:delta-images-equal-classical-boundaries,thm:external-lambda-bockstein,cor:fess-triangle-shift,prop:lambda-rho-ess-only-d0}
  \notready
  Suppose the classical Adams spectral sequence has
  $d_r(x)=y$, with $x\in Z_{r-1}^{s,t}$ and
  $y\in Z_\infty^{s+r,t+r-1}$.  For
  $\delta_n:\nu X/\lambda^n\to\Sigma^{1,-n}\nu X$:
  \begin{enumerate}
    \item if $r\ge n+1$, then
      $d_r^{\delta_n}(x)=\lambda^{r-n-1}y$;
    \item if $r\le n+1$, then
      $d_r^{\delta_n}(\lambda^{n+1-r}x)=y$.
  \end{enumerate}
  The formulas agree when $r=n+1$ and are equalities in the correct ESS
  quotient, whether or not the displayed extension is essential.
\end{proposition}

\begin{proof}
  Start with $n=1$, where the assertion is the $\lambda$-Bockstein theorem.
  Induct on $n$ using the commuting triangle relating multiplication by
  $\lambda$, $\delta_{n-1}$, and $\delta_n$.  Triangle propagation gives the
  formula before quotient indeterminacy.  Division by one $\lambda$ introduces
  only the next classical boundary subgroup; Lemma~\ref{lem:delta-images-equal-classical-boundaries}
  identifies and removes precisely that ESS indeterminacy.
\end{proof}

% SOURCE: aimpaper/main.tex:1068-1131, Proposition prop:6de7d130, finite m case.
% VERIFIED-IN: the printed second case places y in E_infty(nu X); the finite target requires its rho-image.
\begin{proposition}[Finite $\delta_{n,m}$-ESS formula]
  \label{prop:delta-ess-formula-finite}
  \uses{prop:delta-ess-formula-infinite,cor:fess-composition,prop:lambda-rho-ess-only-d0,prop:lambda-quotient-restriction-coherence}
  \notready
  For finite $m>n$, compose the formulas of
  Proposition~\ref{prop:delta-ess-formula-infinite} with the reduction
  $\rho:\nu X\to\nu X/\lambda^{m-n}$.  Thus the displayed target is the image
  of $\lambda^{r-n-1}y$, or of $y$ in the second case, in the finite quotient;
  it is zero once the required $\lambda$ power lies outside that quotient.
  This explicit image corrects the type ambiguity in
  \texttt{aimpaper/main.tex:1088--1094}.
\end{proposition}

\begin{proof}
  Factor $\delta_{n,m}$ as $\delta_{n,\infty}$ followed by $\rho$.  The
  $\rho$-ESS has only its injective-or-zero $d_0$, with no crossing, so ESS
  composition transports the infinite formula to the finite target.  The
  quotient formula decides exactly when the image is zero.
\end{proof}

% SOURCE: aimpaper/main.tex:1138-1144, Corollary cor:2a636737.
\begin{corollary}[$\lambda$-multiple form of the $\delta$ formula]
  \label{cor:delta-ess-lambda-multiples}
  \uses{prop:delta-ess-formula-finite}
  \notready
  In the situation above,
  \[
    d_r^\delta(\lambda^a x)=\lambda^{a+r-n-1}y
  \]
  whenever $0\le a\le n$ and
  $0\le a+r-n-1<m-n$; outside the finite quotient range the right-hand side
  and the extension are zero.
\end{corollary}

\begin{proof}
  Multiply the finite formula by $\lambda^a$, using $\lambda$-linearity of the
  synthetic Adams and extension spectral sequences.  Check the source and
  target quotient ranges directly.
\end{proof}

% SOURCE: aimpaper/main.tex:1146-1150, Remark rmk:qvoewfj.
\begin{theorem}[$\delta$-ESS and classical differential cosets are equivalent]
  \label{thm:delta-classical-equivalence}
  \uses{cor:delta-ess-lambda-multiples,lem:delta-images-equal-classical-boundaries,thm:external-lambda-bockstein}
  \notready
  For every allowed $a,n,m,r$, the equation
  $d_r^\delta(\lambda^a x)=\lambda^{a+r-n-1}y$ holds if and only if the
  classical differential coset is $d_r(x)=y$.  Both sides have the same
  indeterminacy $B_{r-1}$ after applying the stated $\lambda$ power.
\end{theorem}

\begin{proof}
  The forward direction applies the Bockstein comparison to a representative
  and uses the shorter-image lemma to identify the quotient.  The reverse
  direction is the $\lambda$-multiple formula.  Equality of the two boundary
  subgroups makes the correspondence bijective at the level of differential
  cosets, not merely at the level of nonvanishing.
\end{proof}

% SOURCE: aimpaper/main.tex:1152-1163, Example 4.13.
\begin{proposition}[$\delta_1,\delta_2,\delta_3$ regression]
  \label{prop:delta-14-stem-regression}
  \uses{prop:synthetic-14-stem-regression,cor:delta-ess-lambda-multiples}
  \notready
  The $d_2(h_4)=h_0h_3^2$ and
  $d_3(h_0h_4)=h_0d_0$ examples yield exactly the eight
  $\delta_1,\delta_2,\delta_3$ extensions and $\lambda$ powers listed in
  Example 4.13: two for $\delta_1$, three for $\delta_2$, and three for
  $\delta_3$.
\end{proposition}

\begin{proof}
  Substitute $(n,m)=(1,\infty),(2,\infty),(3,\infty)$ and each admissible
  value of $a$ into Corollary~\ref{cor:delta-ess-lambda-multiples}.  The
  synthetic $E_\infty$ regression verifies that every source and target exists
  in the displayed tridegree.
\end{proof}

% SOURCE: aimpaper/main.tex:1165-1173, Definition def:classicalcrossdiff.
\begin{definition}[Classical crossing on an $E_{n+1}$ page]
  \label{def:classical-differential-crossing}
  \uses{def:classical-adams-ss,def:ss-permanent-einfty-class}
  \notready
  Let $r\ge n+1$ and $d_r(x)=y$, with $x$ in filtration $s$.  A crossing on
  the $E_{n+1}$ page is an essential differential
  $d_{r-a-b}(x')=y'$ where the source filtration is $s+a$, the target
  filtration is $s+r-b$, $0<a\le n-1$, and
  $0\le b\le r-n-1$.
\end{definition}

% SOURCE: aimpaper/main.tex:1218-1234, Proposition prop:cross-dr-En.
\begin{proposition}[Classical crossing iff $\delta_n$ crossing]
  \label{prop:classical-crossing-iff-delta-crossing}
  \uses{def:classical-differential-crossing,def:fess-crossing,thm:delta-classical-equivalence,thm:external-synthetic-einfty-nu,thm:external-synthetic-einfty-quotient,prop:synthetic-einfty-first-quotient}
  \notready
  A classical $d_r(x)=y$ has a crossing on $E_{n+1}$ if and only if the
  corresponding extension
  $d_r^{\delta_n}(x)=\lambda^{r-n-1}y$ has a crossing in the synthetic
  $\delta_n$-ESS.
\end{proposition}

\begin{proof}
  A synthetic crossing has the unique degree form
  $d_{r-a-b}^{\delta_n}(\lambda^a x')=\lambda^{r-n-1-b}y'$.  The two
  sides are both empty when $n=1$: the classical range would require
  $0<a\le0$, while mod-$\lambda$ collapse and unique detection exclude an
  additional synthetic representative.  For $n\ge2$, the two $E_\infty$
  formulas force exactly the inequalities in
  Definition~\ref{def:classical-differential-crossing}.  Apply the
  $\delta$--classical equivalence in each direction, preserving essentiality
  through the shared boundary quotient.  The first-quotient proposition is
  used only in the $n=1$ branch.
\end{proof}

% SOURCE: aimpaper/main.tex:1236-1238, Remark nocrossE2.
\begin{corollary}[No classical crossing on $E_2$]
  \label{cor:classical-e2-no-crossing}
  \uses{def:classical-differential-crossing,prop:classical-crossing-iff-delta-crossing}
  \notready
  No classical Adams differential has a crossing on the $E_2$ page, and every
  corresponding $\delta_1$ extension has no crossing.
\end{corollary}

\begin{proof}
  Setting $n=1$ would require an integer $a$ with $0<a\le0$.  Hence no
  classical crossing exists; transfer the conclusion through
  Proposition~\ref{prop:classical-crossing-iff-delta-crossing}.
\end{proof}

% SOURCE: aimpaper/main.tex:1240-1271, Example exam:classcrossdiff.
\begin{proposition}[Stem-$38$ crossing regression]
  \label{prop:stem38-crossing-regression}
  \uses{prop:classical-crossing-iff-delta-crossing,cor:classical-e2-no-crossing,def:external-evidence}
  \notready
  Given the typed computation evidence
  $d_3(e_1)=h_1t$ and
  $d_4(h_0h_3h_5)=h_0^2x$, the former is a crossing of the latter on the
  $E_3$ and $E_4$ pages, but not on $E_2$ or on pages $E_r$ for $r\ge5$.
  The associated $\delta_n$ extensions and their weights are exactly those in
  Example 4.18.
\end{proposition}

\begin{proof}
  Check the four inequalities in the crossing definition for $n=1,2,3$ and
  $n\ge4$, then transport the two computed differential cosets using the
  crossing equivalence.  The synthetic $E_\infty$ formulas verify which
  $\lambda$ multiples remain represented.
\end{proof}

\section{Classical page-extension theorems}

% SOURCE: aimpaper/main.tex:1541-1564 and 2084-2099, the finite Ext-product
% enumeration used both for no-crossing and for page stretching.  This node is
% deliberately stated without assuming that the E_infinity extension exists.
\begin{proposition}[Hopf crossing-obstruction exclusion]
  \label{evidence:hopf-crossing-obstruction-exclusion}
  \uses{def:page-extension-crossing,prop:page-extension-loss-certificate,thm:external-adams-one-line,def:ss-elementwise-page-presentation,def:ss-cycle-boundary-survival,def:external-evidence,evidence:normalized-hopf-detection}
  \notready
  For the candidate Hopf relation of length two from
  $h_0h_4^2$ to $h_0p$, a finite evidence record enumerates every
  degree-allowed crossing tuple and every tuple satisfying the
  page-loss-certificate bounds between $E_3$ and any finite page.  The list is
  empty.  A separate part enumerates every permanent source and target in the
  same forced bidegrees and excludes every actual shorter
  $(\nu_{\mathrm{Hopf}},E_\infty)$-extension; this is not inferred from the
  finite list.  Both the finite-page reductions and this permanent-candidate
  enumeration use the Hopf component of the single stored normalized
  triangle lift from Proposition~\ref{evidence:normalized-hopf-detection},
  together with its retained equality
  $\widehat{\nu_{\mathrm{Hopf}}}=[h_2]$.  Its decisive typed product check is
  \[
    h_0^2h_4^2h_2=0\ne h_0p.
  \]
  The same record supplies the independent typing facts
  \[
    h_0h_4^2\in Z_\infty(S^3),
    \qquad h_0p\in Z_\infty(S^0).
  \]
  The first follows from the known $d_2(h_5)=h_0h_4^2$ and
  $B_2\subseteq Z_\infty$; the second is a complete outgoing-differential
  exclusion in the target bidegree.  Thus both the source and target remain
  typed on every later page.  It also certifies
  $h_0p\notin B_2$ and that no shorter Hopf-ESS image reaches its complete
  target coset.  Consequently, whenever the candidate length-two relation
  exists at a finite or infinite level, it is essential.
  Thus the record contains both a finite-page obstruction list and an
  independently computed untruncated candidate list.  Neither branch assumes
  nor concludes that the candidate relation already exists at its respective
  finite or infinite level.
\end{proposition}

\begin{proof}
  Enumerate the finitely many integers $a,a',b$ allowed by the crossing and
  loss-certificate inequalities and then enumerate the Ext groups in their
  forced source and target bidegrees.  Intersect the same finite bidegree lists
  with $Z_\infty$ to obtain the independent untruncated-candidate list.  The
  stored equality $\widehat{\nu_{\mathrm{Hopf}}}=[h_2]$ identifies, on every
  finite reduction and on the untruncated list, the only possible shorter Hopf
  product with the displayed zero class, whereas the proposed target is
  nonzero.  Store
  the degree bounds, cycle-range checks, complete incoming/outgoing target
  search, and typed product inequality in the evidence record.  The source cycle fact is
  also obtained from the one-line differential and the inclusion of boundaries
  into permanent representatives.  No inverse-limit representative or
  $E_\infty$ extension is used in this enumeration.
\end{proof}

% SOURCE: aimpaper/main.tex:1284-1298, E_0 formulas before Definition def:6c076a33.
\begin{proposition}[$E_0$ presentation for the reduced-map ESS]
  \label{prop:hat-map-ess-e0}
  \uses{def:hat-map-page-reduction,def:synthetic-extension-ss,thm:external-synthetic-einfty-nu,thm:external-synthetic-einfty-quotient,prop:synthetic-einfty-first-quotient}
  \notready
  For finite $r$, in weight $t-k+e(f)$, the $E_0$ page of the
  $\widehat f_{r-1}$-ESS is
  \[
  \frac{Z_{r-1-k}^{s,t}(X)}{B_{1+k}^{s,t}(X)}\oplus
  \frac{Z_{r-1-k+e(f)}^{s,t}(Y)}{B_{1+k-e(f)}^{s,t}(Y)}.
  \]
  This displayed formula is for finite $r$.  At $r=\infty$, the $E_0$ page
  of the $\widehat f$-ESS is instead
  \[
  \frac{Z_\infty^{s,t}(X)}{B_{1+k}^{s,t}(X)}\oplus
  \frac{Z_\infty^{s,t}(Y)}{B_{1+k-e(f)}^{s,t}(Y)};
  \]
  no expression such as $\infty-1-k$ is formed.
  For finite $r$, its length-$n$ differential is a map from
  \[
    {}^{\widehat f_{r-1}}\!Z_{n-1}^{s,t,t-k}(X)
      \subseteq Z_{r-1-k}^{s,t}(X)/B_{1+k}^{s,t}(X)
  \]
  to
  \[
    \frac{Z_{r-1-k-n+e(f)}^{s+n,t+n}(Y)/
      B_{1+k+n-e(f)}^{s+n,t+n}(Y)}
    {{}^{\widehat f_{r-1}}\!B_{n-1}^{s+n,t+n,t-k}(Y)}.
  \]
  It vanishes for degree reasons when $n<\operatorname{AF}(f)$ or
  $n>r-2-k+e(f)$.
  At $r=\infty$, replace $\widehat f_{r-1}$ by $\widehat f$ and every cycle
  numerator in these source and target groups by $Z_\infty$; only the lower
  vanishing bound $n<\operatorname{AF}(f)$ remains.
\end{proposition}

\begin{proof}
  If $r=2$, substitute the source and target weights into the first-quotient
  edge presentation.  If $3\le r<\infty$, use the finite-quotient
  $E_\infty$ formula at exponent $r-1\ge2$.  If $r=\infty$, use the
  untruncated $Z_\infty/B_j$ formula separately.  Then apply the elementwise ESS presentation,
  simplify the two cycle and boundary subscripts over integers, and check the
  quotient range before converting to a finite page index.
\end{proof}

% SOURCE: aimpaper/main.tex:1391-1402, Example exam:extonEn(1).
\begin{proposition}[Hopf-map $E_2$ page extension]
  \label{prop:hopf-e2-page-extension}
  \uses{def:page-extension,evidence:normalized-hopf-detection,prop:synthetic-einfty-first-quotient,def:external-evidence}
  \notready
  Let $\nu_{\mathrm{Hopf}}:S^3\to S^0$ denote the classical Hopf map.  Typed Ext
  product evidence supplies the essential extension
  \[
    d_1^{\nu_{\mathrm{Hopf}},E_2}(h_5)=h_2h_5.
  \]
  The notation $\nu_{\mathrm{Hopf}}$ is kept distinct from the synthetic functor
  $\nu$.
\end{proposition}

\begin{proof}
  Reduce the precise lifted-triangle component identified in
  Proposition~\ref{evidence:normalized-hopf-detection} modulo $\lambda$.
  The first-quotient uniqueness statement identifies its reduction with
  multiplication by $h_2$, so the stated extension is exactly the supplied
  Ext product.  Thus this node and the later $E_3/E_\infty$ nodes use the same
  chosen normalized Hopf map.
\end{proof}

% SOURCE: aimpaper/main.tex:1403-1423, Example exam:extonEn(2); proved later by Example exam:Mahowald.
\begin{proposition}[Hopf-map $E_3$ page extension]
  \label{prop:hopf-e3-page-extension}
  \uses{def:page-extension,thm:generalized-mahowald,prop:mahowald-cofiber-regression,evidence:hopf-crossing-obstruction-exclusion}
  \notready
  The normalized Hopf map has the essential extension
  \[
    d_2^{\nu_{\mathrm{Hopf}},E_3}(h_0h_4^2)=h_0p.
  \]
  Equivalently, modulo $\lambda^2$ there are detected representatives with
  $[h_0h_4^2][h_2]=[\lambda h_0p]$.
\end{proposition}

\begin{proof}
  This is the conclusion of the cofiber regression obtained from the
  generalized Mahowald trick.  Transport that synthetic extension through
  Definition~\ref{def:page-extension} and verify the displayed bidegrees.  The
  independent Hopf obstruction certificate says that $h_0p$ is nonzero modulo
  $B_2$ and is not a shorter Hopf-ESS image, so the complete target coset does
  not contain zero; hence the extension is essential.
\end{proof}

% SOURCE: aimpaper/main.tex:1424-1433, Example exam:extonEn(3); proved later by Example exam:stretchext.
\begin{proposition}[Hopf-map $E_\infty$ page extension]
  \label{prop:hopf-einfty-page-extension}
  \uses{prop:hopf-e3-page-extension,prop:page-stretch-regression}
  \notready
  The preceding relation lifts to an essential extension
  \[
    d_2^{\nu_{\mathrm{Hopf}},E_\infty}(h_0h_4^2)=h_0p.
  \]
\end{proposition}

\begin{proof}
  This is the essential untruncated relation proved by the page-stretch
  regression.  Its coherent inverse-limit construction supplies the relation,
  and the independent nonzero target-coset clause supplies essentiality.
\end{proof}

% SOURCE: aimpaper/main.tex:1391-1433, Example exam:extonEn(4).
\begin{proposition}[Multiplication-by-two page-extension regression]
  \label{prop:two-page-extension-regression}
  \uses{def:page-extension-coset,prop:no-crossing-iff-uniform-detection,def:external-evidence}
  \notready
  Given the stated $14$-stem evidence and the comparison
  $\lambda[h_0]=2$, multiplication by two has
  $d_2^{2,E_\infty}(h_0h_3^2)=0$, while
  $d_1^{2,E_\infty}(d_0)=h_0d_0$ is essential and makes the alternative target
  $d_2^{2,E_\infty}(h_0h_3^2)=h_0d_0$ inessential.
\end{proposition}

\begin{proof}
  Evaluate multiplication by $[h_0]$ on the supplied synthetic homotopy group.
  Its only nonzero leading target comes from the higher-filtration class
  $[\lambda d_0]$.  The representative and inessentiality criteria yield the
  three asserted page-extension statuses.
\end{proof}

% SOURCE: aimpaper/main.tex:1515-1539, Proposition prop:cross-f-Er.
\begin{proposition}[Page crossing iff synthetic crossing]
  \label{prop:page-crossing-iff-synthetic-crossing}
  \uses{def:page-extension-crossing,def:fess-crossing,prop:synthetic-einfty-map-compatibility,prop:synthetic-ess-weightwise-transport}
  \notready
  An $(f,E_r)$-extension has a page crossing if and only if its defining
  $\widehat f_{r-1}$-ESS extension has a crossing for finite $r$.  For
  $r=\infty$, the equivalent objects are the actual untruncated
  $(f,E_\infty)$ crossing and an actual crossing in the $\widehat f$-ESS.
\end{proposition}

\begin{proof}
  A synthetic crossing necessarily has the form
  $d_{n-a-b}^{\widehat f_{r-1}}(\lambda^ax')
   =\lambda^{n-b-e(f)}y'$.  Compare the quotient of exponent $r-1$ with that
  of exponent $r-1-a$.  Dividing by $\lambda^a$ introduces exactly
  $B_{1+n-b-e(f)}$; therefore essentiality is equivalent to the nonboundary
  condition in the finite page-crossing definition.

  At $r=\infty$, start with an actual untruncated shorter page extension on
  permanent representatives.  The infinite-endpoint clause of
  Proposition~\ref{prop:synthetic-einfty-map-compatibility} identifies
  multiplication by $\lambda^a$ with the corresponding boundary-cutoff
  quotient on $Z_\infty/B_j$.  This converts that actual page witness to an
  actual crossing in the untruncated $\widehat f$-ESS, and the same quotient
  calculation gives the converse.  No finite solution tower is promoted to an
  infinite witness in this argument.
\end{proof}

% SOURCE: aimpaper/main.tex:1541-1579, Example exam:Ercross.
\begin{proposition}[Page-crossing regression]
  \label{prop:page-crossing-regression}
  \uses{prop:hopf-e2-page-extension,prop:hopf-e3-page-extension,prop:hopf-einfty-page-extension,evidence:hopf-crossing-obstruction-exclusion,prop:two-page-extension-regression,prop:page-crossing-iff-synthetic-crossing,def:external-evidence}
  \notready
  The four examples of Section 5 have the stated crossing behavior: the
  $E_2$, $E_3$, and $E_\infty$ Hopf-map extensions have no crossings, while
  the zero multiplication-by-two extension is crossed by the essential
  length-one $d_0$ target.  The no-crossing certificates include the Ext
  relation $h_0^2h_4^2h_2=0\ne h_0p$.
\end{proposition}

\begin{proof}
  There is no positive $a$ in the $E_2$ crossing range.  For the actual
  Hopf-map $E_3$ and $E_\infty$ extensions, apply the independent finite
  obstruction certificate.  Its separate permanent-candidate enumeration
  excludes every actual untruncated shorter extension without having assumed
  the candidate $E_\infty$ relation.  For multiplication by two, the essential
  $d_1$ extension from $d_0$ meets the crossing inequalities and witnesses the
  failure of no-crossing.
\end{proof}

% SOURCE: aimpaper/main.tex:1581-1597, unlabelled Proposition 5.14.
\begin{proposition}[$E_\infty$ page extension gives a classical extension]
  \label{prop:einfty-page-to-classical-extension}
  \uses{def:page-extension,thm:external-lambda-inversion,def:reindexed-ss-morphism,def:f-extension-essential}
  \notready
  If $x\in Z_\infty^{s,t}(X)$ and
  $y\in Z_\infty^{s+n,t+n}(Y)$ satisfy
  $d_n^{f,E_\infty}(x)=y$, then after passing to classical $E_\infty$ cosets,
  \[
    d_n^f(x+B_\infty)=y+B_\infty.
  \]
  The resulting classical extension may be inessential.
\end{proposition}

\begin{proof}
  Invert $\lambda$ in the $\widehat f$-ESS.  Generic-fiber equivalence
  identifies its source and target with the classical $f$-ESS and sends the
  normalized map to $f$.  Naturality of the induced reindexed page maps sends
  the synthetic differential coset to the displayed classical coset.
\end{proof}

\section{Generalized Leibniz and Mahowald rules}

% SOURCE: aimpaper/main.tex:1606-1652, Theorem thm:e73f481e.
\begin{theorem}[Generalized Leibniz rule]
  \label{thm:generalized-leibniz}
  \uses{prop:delta-ess-formula-infinite,thm:delta-classical-equivalence,prop:classical-crossing-iff-delta-crossing,def:page-extension,prop:page-crossing-iff-synthetic-crossing,cor:fess-square-naturality}
  \notready
  Let $f:X\to Y$, $2\le n\le r$,
  $e(f)\le m\le n-2+e(f)$, and $l\ge e(f)$.  Suppose
  \[
  \begin{aligned}
    &x\in Z_{r-1}^{s,t}(X),&&
    y\in Z_{r-1-m+e(f)}^{s+m,t+m}(Y),\\
    &x_\infty\in Z_\infty^{s+r,t+r-1}(X),&&
    y_\infty\in Z_\infty^{s+r+l,t+r+l-1}(Y),
  \end{aligned}
  \]
  and assume:
  \begin{enumerate}
    \item $d_r(x)=x_\infty$;
    \item $d_m^{f,E_n}(x)=y$;
    \item $d_l^{f,E_\infty}(x_\infty)=y_\infty$;
    \item either the first differential has no crossing on $E_n$, or the
      second page extension has no crossing;
    \item the third page extension has no crossing.
  \end{enumerate}
  Then $d_{r+l-m}(y)=y_\infty$.
\end{theorem}

\begin{proof}
  Form the square whose horizontal maps are
  $\widehat f_{n-1}$ and $\widehat f$ and whose vertical maps are the two
  connecting maps from reduction modulo $\lambda^{n-1}$.  Translate (1) by the
  $\delta$ formula and (2)--(3) by the definition of page extension.  The two
  crossing comparison theorems turn hypothesis (4) into synchronized-source
  no-crossing and hypothesis (5) into target no-crossing.  Apply ESS square
  propagation.  Finally use the $\delta$--classical equivalence to remove the
  displayed $\lambda$ powers and obtain the asserted classical differential
  coset.
\end{proof}

% SOURCE: aimpaper/main.tex:1664-1692, Example exam:Leibnizyes.
\begin{proposition}[Positive Leibniz regression]
  \label{prop:leibniz-positive-regression}
  \uses{thm:generalized-leibniz,thm:external-adams-one-line,prop:hopf-e2-page-extension,prop:hopf-einfty-page-extension,cor:classical-e2-no-crossing,prop:page-crossing-regression}
  \notready
  The inputs
  $d_2(h_5)=h_0h_4^2$,
  $d_1^{\nu_{\mathrm{Hopf}},E_2}(h_5)=h_2h_5$, and
  $d_2^{\nu_{\mathrm{Hopf}},E_\infty}(h_0h_4^2)=h_0p$
  satisfy the generalized Leibniz hypotheses and give
  \[
    d_3(h_2h_5)=h_0p.
  \]
\end{proposition}

\begin{proof}
  Substitute $(n,r,m,l,e(f))=(2,2,1,2,1)$.  The $E_2$ classical crossing
  condition is automatic, and the two Hopf-map page extensions have the
  no-crossing certificates from the regression.  The theorem gives a
  differential of length $2+2-1=3$ with the claimed target.
\end{proof}

% SOURCE: aimpaper/main.tex:1694-1720, Example exam:Leibnizno.
\begin{proposition}[No-crossing is necessary]
  \label{prop:leibniz-negative-regression}
  \uses{thm:generalized-leibniz,prop:two-page-extension-regression,prop:page-crossing-regression,def:external-evidence}
  \notready
  For $f=2$ and the computed differential
  $d_2(h_4)=h_0h_3^2$, all generalized Leibniz hypotheses except target
  no-crossing can be met using
  $d_2^{2,E_\infty}(h_0h_3^2)=0$.  The missing hypothesis is false because
  $d_1^{2,E_\infty}(d_0)=h_0d_0$ crosses it.  Omitting that hypothesis would
  assert the false zero differential $d_3(h_0h_4)=0$.
\end{proposition}

\begin{proof}
  Evaluate every hypothesis using the page-extension regression and the typed
  classical computation.  The length-one extension meets the crossing bounds,
  so the generalized theorem is inapplicable.  The known nonzero differential
  on $h_0h_4$ supplies the regression check that the weakened statement is
  false.
\end{proof}

% SOURCE: aimpaper/main.tex:1733-1753, Remark rem:chuaerror; Chua Definition 12.5 and Theorem 12.9.
\begin{proposition}[Counterexample to the maximal-extension rule]
  \label{prop:chua-rule-counterexample}
  \uses{prop:leibniz-negative-regression,def:external-result,def:external-evidence}
  \notready
  The maximal-$[h_0]$ lift of $h_0h_3^2$ in the stated synthetic $14$-stem
  example maps to zero, while another lift maps to $[h_0d_0]$.  Consequently
  the rule quoted as Chua Theorem 12.9, when read without a no-crossing
  hypothesis, predicts the false conclusion of
  Proposition~\ref{prop:leibniz-negative-regression}.
\end{proposition}

\begin{proof}
  Use the two distinct detected lifts and their products from the supplied
  evidence.  The zero product is maximally $\lambda$-divisible under the cited
  definition.  Substitution into the quoted rule yields the false zero target,
  so these data form a concrete counterexample; no claim about a corrected
  version is made.
\end{proof}

% SOURCE: aimpaper/main.tex:1780-1917, complete hypotheses and degree package for Theorem thm:158d451a.
\begin{definition}[Generalized Mahowald input package]
  \label{def:mahowald-input-data}
  \uses{prop:normalized-synthetic-triangle,def:page-extension,def:classical-differential-crossing}
  \notready
  A Mahowald input package consists of a distinguished triangle
  $X\xrightarrow fY\xrightarrow gZ\xrightarrow h\Sigma X$ with
  $e(f)+e(g)+e(h)=1$, integers $s,t,n,m,l$, and
  \[
    r=n+m+l,\quad n_1=n-e(f)\ge1,\quad
    m_1=m-e(g)\ge0,\quad l_1=l-e(h)\ge0,\quad
    r'=r-m_1=n_1+l_1+1.
  \]
  It contains classes
  \[
  \begin{aligned}
    x&\in Z_{n_1}^{s+l,t+l-1}(X),&
    y&\in Z_{m_1+1}^{s+n+l,t+n+l-1}(Y),\\
    \bar x&\in Z_{r-1}^{s,t}(Z),&
    \bar y&\in Z_\infty^{s+r,t+r-1}(Z),
  \end{aligned}
  \]
  proofs of $d_l^{h,E_{r'}}(\bar x)=x$, $d_r(\bar x)=\bar y$,
  $d_m^{g,E_{m_1+2}}(y)=\bar y$, and a proof that either the first page
  extension has no crossing or the classical differential has no crossing on
  $E_{r'}$.
\end{definition}

% SOURCE: aimpaper/main.tex:1780-1917, representative choices in the proof of Theorem thm:158d451a.
\begin{lemma}[Synchronized Mahowald representatives]
  \label{lem:mahowald-synchronized-representatives}
  \uses{def:mahowald-input-data,prop:no-crossing-iff-uniform-detection,prop:classical-crossing-iff-delta-crossing,prop:page-crossing-iff-synthetic-crossing,thm:delta-classical-equivalence}
  \notready
  For a Mahowald input package, one can choose
  \[
    [\bar x]\in\pi_{t-s,t}(\nu Z/\lambda^{r'-1})
  \]
  detected by $\bar x$ such that
  $\widehat h_{r'-1}[\bar x]$ is detected by
  $\lambda^{l_1}x$ in
  $\pi_{t-s-1,t-1+e(h)}(\nu X/\lambda^{r'-1})$ and
  $\delta_{r'-1}[\bar x]$ is detected by
  $\lambda^{m_1}\bar y$ in
  $\pi_{t-s-1,t+r'-1}(\nu Z/\lambda^{m_1+1})$.
\end{lemma}

\begin{proof}
  Translate the page extension and classical differential into the two
  synthetic ESS extensions in the common quotient.  The assumed no-crossing
  of at least one extension invokes uniform detection for every representative
  chosen to realize the other.  Hence one representative realizes both
  leading terms simultaneously.
\end{proof}

% SOURCE: aimpaper/main.tex:1886-1907, the representative for y, degree-uniqueness, and the application of May's pullback lemma.
\begin{lemma}[May-compatible lift in the Mahowald diagram]
  \label{lem:mahowald-compatible-quotient-lift}
  \uses{def:mahowald-input-data,lem:mahowald-synchronized-representatives,thm:external-may-smash-boundary,prop:synthetic-ess-weightwise-transport,prop:fextension-iff-detected,thm:external-synthetic-einfty-quotient,prop:synthetic-einfty-first-quotient,prop:normalized-synthetic-triangle,def:lambda-rho-delta-triangle}
  \notready
  For a Mahowald input package, choose the synchronized representative
  $[\bar x]$ of Lemma~\ref{lem:mahowald-synchronized-representatives}.
  There are a representative

  \[
    [y]\in\pi_{t-s-1,t+n+l-1}(\nu Y/\lambda^{m_1+1})
  \]
  detected by $y$ and a class

  \[
    u\in\pi_{t-s-1,t-1+e(h)}(\nu X/\lambda^r)
  \]
  detected by $\lambda^{l_1}x$ such that

  \[
    \rho(u)=\widehat h_{r'-1}([\bar x]),
    \qquad
    \widehat f_r(u)=\lambda^{r'-1}[y].
  \]
\end{lemma}

\begin{proof}
  The fourth page-extension hypothesis and the representative criterion choose
  $[y]$ whose image under $\widehat g_{m_1+1}$ is detected by
  $\lambda^{m_1}\bar y$.  If $m_1=0$, the first-quotient collapse and its
  unique-detection conclusion make that mod-$\lambda$ homotopy class unique;
  if $m_1\ge1$, the finite-quotient $E_\infty$ formula gives uniqueness in the
  indicated top weight.  Hence in both cases
  $\widehat g_{m_1+1}([y])$ equals
  $\delta_{r'-1}([\bar x])$.  Apply May's smash-boundary theorem to the
  normalized synthetic triangle and the $\lambda$--$\rho$--$\delta$ triangle.
  Its pullback class is $u$, and its two structure maps give exactly the two
  displayed equalities.
\end{proof}

% SOURCE: aimpaper/main.tex:1900-1910, equation eq:f4848f29 and the ensuing one-way survival deduction.
\begin{lemma}[A compatible quotient lift gives source survival]
  \label{lem:mahowald-quotient-lift-survival}
  \uses{def:mahowald-input-data,lem:mahowald-synchronized-representatives,thm:external-synthetic-einfty-quotient,prop:synthetic-einfty-map-compatibility}
  \notready
  For a Mahowald input package, suppose

  \[
    u\in\pi_{t-s-1,t-1+e(h)}(\nu X/\lambda^r)
  \]
  is detected by $\lambda^{l_1}x$ and its $\rho$-image is the synchronized
  representative detected by $\lambda^{l_1}x$ modulo $\lambda^{r'-1}$.
  Then

  \[
    x\in Z_{n+m+e(h)}^{s+l,t+l-1}(X).
  \]
  No converse for a preselected representative in the smaller quotient is
  asserted.
\end{lemma}

\begin{proof}
  Apply the finite-quotient $E_\infty$ formula to the class $u$.  Its exponent
  after removing the displayed $\lambda^{l_1}$ is
  $r-l_1=n+m+e(h)$, so the existence of $u$ places $x$ in the corresponding
  cycle subgroup.  Compatibility of $\rho$ verifies that its reduction is the
  synchronized smaller-quotient representative; it is not used to manufacture
  a lift from survival alone.
\end{proof}

% SOURCE: aimpaper/main.tex:1903-1917, final lambda-division step of Theorem thm:158d451a.
\begin{lemma}[$\lambda$-division indeterminacy in the Mahowald diagram]
  \label{lem:mahowald-lambda-division}
  \uses{def:mahowald-input-data,prop:synthetic-einfty-map-compatibility,def:page-extension-coset}
  \notready
  For a Mahowald input package, if the smash diagram gives
  \[
    d_n^{\widehat f_r}(\lambda^{l_1}x)=\lambda^{r'-1}y,
  \]
  then division by $\lambda^{l_1}$ is defined in the quotient of exponent
  $r-l_1=n+m+e(h)$ and changes the target exactly by
  $B_{r'}^{s+n+l,t+n+l-1}(Y)$.  The result is
  \[
    d_n^{f,E_{n+m+1+e(h)}}(x)
       \equiv y\pmod{B_{r'}^{s+n+l,t+n+l-1}(Y)}.
  \]
\end{lemma}

\begin{proof}
  Apply the quotient-map description of multiplication by $\lambda^{l_1}$.
  Its kernel is the boundary cutoff enlarged by $l_1$; after substituting the
  definitions of $n_1,m_1,l_1,r'$, this cutoff is $B_{r'}$.  Compare with the
  complete page-extension target coset to obtain the asserted equivalence.
\end{proof}

% SOURCE: aimpaper/main.tex:1780-1917, Theorem thm:158d451a.
\begin{theorem}[Generalized Mahowald trick]
  \label{thm:generalized-mahowald}
  \uses{def:mahowald-input-data,prop:normalized-synthetic-triangle,def:lambda-rho-delta-triangle,lem:mahowald-synchronized-representatives,lem:mahowald-compatible-quotient-lift,lem:mahowald-quotient-lift-survival,lem:mahowald-lambda-division}
  \notready
  Let
  $X\xrightarrow fY\xrightarrow gZ\xrightarrow h\Sigma X$ be distinguished
  with $e(f)+e(g)+e(h)=1$.  Put
  $r=n+m+l$, $n_1=n-e(f)\ge1$,
  $m_1=m-e(g)\ge0$, $l_1=l-e(h)\ge0$, and
  $r'=r-m_1=n_1+l_1+1$.  Suppose
  \[
  \begin{aligned}
    x&\in Z_{n_1}^{s+l,t+l-1}(X),&
    y&\in Z_{m_1+1}^{s+n+l,t+n+l-1}(Y),\\
    \bar x&\in Z_{r-1}^{s,t}(Z),&
    \bar y&\in Z_\infty^{s+r,t+r-1}(Z),
  \end{aligned}
  \]
  and assume
  \begin{enumerate}
    \item $d_l^{h,E_{r'}}(\bar x)=x$;
    \item $d_r(\bar x)=\bar y$;
    \item either (1) has no page crossing or (2) has no crossing on $E_{r'}$;
    \item $d_m^{g,E_{m_1+2}}(y)=\bar y$.
  \end{enumerate}
  Then $x\in Z_{n+m+e(h)}$ and
  \[
    d_n^{f,E_{n+m+1+e(h)}}(x)
       \equiv y\quad\text{modulo }B_{r'}.
  \]
\end{theorem}

\begin{proof}
  Smash the normalized synthetic triangle with the quotient triangle of
  exponents $r$ and $r'-1$.  The synchronization lemma chooses the common
  representative of $\bar x$.  The May-compatible-lift lemma uses condition
  (4), the degree-unique common target, and May's pullback theorem to produce a
  single class mapping to both sides of the diagram.  The one-way
  quotient-lift lemma turns that actual $\rho$ lift into the stronger survival
  statement for $x$.  Its other image is the normalized $f$ relation; divide
  by $\lambda^{l_1}$ using the final lemma to obtain the page extension with
  exactly $B_{r'}$ indeterminacy.
\end{proof}

% SOURCE: aimpaper/main.tex:1930-1977, Example exam:Mahowald; computed cofiber d_2 is external evidence.
\begin{proposition}[Cofiber regression for the Mahowald trick]
  \label{prop:mahowald-cofiber-regression}
  \uses{thm:generalized-mahowald,lem:normalized-hopf-exactness-case,def:hopf-cofiber-cell-notation,cor:classical-e2-no-crossing,def:external-evidence}
  \notready
  Given the computed differential
  $d_2(h_0h_4^2[4])=h_0p[0]$ in the Adams spectral sequence of
  $S^0/\nu_{\mathrm{Hopf}}$ and the two displayed Ext edge-map facts, the
  generalized Mahowald trick proves that $h_0h_4^2$ survives to $E_3$ and
  \[
    d_2^{\nu_{\mathrm{Hopf}},E_3}(h_0h_4^2)=h_0p.
  \]
\end{proposition}

\begin{proof}
  Use Lemma~\ref{lem:normalized-hopf-exactness-case} to instantiate the Hopf
  triangle with $(e(f),e(g),e(h))=(1,0,0)$ and its one chosen lifted triangle.
  Substitute $m=l=0$, $n=r=r'=2$, $n_1=1$, and $m_1=l_1=0$ into the theorem.
  Bottom- and top-cell Ext exactness give the two length-zero page extensions;
  the supplied cofiber differential gives condition (2), and $E_2$
  no-crossing gives condition (3).  The theorem produces the stated
  extension with zero additional indeterminacy in this degree.
\end{proof}

% SOURCE: aimpaper/main.tex:1981-2068, Proposition prop:dec738d3, first part.
\begin{proposition}[Restriction of a page extension]
  \label{prop:page-extension-restrict}
  \uses{def:page-extension,cor:fess-map-of-stable-pages,prop:lambda-rho-ess-only-d0,prop:lambda-quotient-restriction-coherence}
  \notready
  If $d_n^{f,E_r}(x)=y$ and $2\le r'<r$, reduction along
  \[
    \rho_X:\nu X/\lambda^{r-1}\to\nu X/\lambda^{r'-1},
    \qquad
    \rho_Y:\nu Y/\lambda^{r-1}\to\nu Y/\lambda^{r'-1}
  \]
  gives
  $d_n^{f,E_{r'}}(x)=y$ whenever the source and target are defined.
\end{proposition}

\begin{proof}
  The square of reduced normalized maps commutes.  Since both vertical
  $\rho$-ESS have only $d_0$, the stable-page map theorem induces a morphism of
  the two $\widehat f$ extension spectral sequences.  Naturality sends the
  length-$n$ differential to its earlier-page version.
\end{proof}

% PROOF-SOURCE: aimpaper/main.tex:1981-2068, proof of Proposition
% prop:dec738d3; project exact-obstruction reformulation of the representative
% lifting calculation needed for the later-page converse.
\begin{lemma}[Exact obstruction for finite page-extension restriction]
  \label{lem:page-extension-restriction-exact-obstruction}
  \uses{def:page-extension,def:page-extension-coset,def:page-extension-crossing,def:synthetic-ess-solution-tower,prop:lambda-quotient-restriction-coherence,cor:fess-map-of-stable-pages,prop:lambda-rho-ess-only-d0,prop:fextension-inessential-iff-higher}
  \notready
  Let $2\le r'<r<\infty$, let
  $e(f)\le n\le r'-2+e(f)$, and suppose
  \[
    x\in Z_{r-1}^{s,t}(X),\qquad
    y\in Z_{r-1-n+e(f)}^{s+n,t+n}(Y).
  \]
  Write $\operatorname{Sol}_{r}^{n}(x,y)$ and
  $\operatorname{Sol}_{r'}^{n}(x,y)$ for the representative solution fibers
  of Definition~\ref{def:synthetic-ess-solution-tower} for
  $\widehat f_{r-1}$ and $\widehat f_{r'-1}$.  Reduction defines a restriction
  map and a pointed obstruction sequence
  \[
    \operatorname{Sol}_{r}^{n}(x,y)
      \xrightarrow{\ \rho_*\ }
    \operatorname{Sol}_{r'}^{n}(x,y)
      \xrightarrow{\ \operatorname{ob}_{r,r'}\ }
    \operatorname{Ob}_{r,r'}^{n}(x,y).
  \]
  Here $\operatorname{Ob}_{r,r'}^{n}(x,y)$ is the cokernel of restriction on
  the groups of representative equations after the target is divided by its
  complete shorter-image coset.  The sequence is exact in the pointed sense:
  for every earlier-page solution $u$, the equality
  \[
    \operatorname{ob}_{r,r'}(u)=0
  \]
  holds exactly when $u\in\operatorname{im}(\rho_*)$.

  Every nonzero obstruction has a representative with
  $0<a'\le n-e(f)$, $0\le b\le n-a'-e(f)$, and classes
  \[
    x'\in Z_{r'-1-a'}^{s+a',t+a'}(X)
       \setminus Z_{r-1-a'}^{s+a',t+a'}(X),
    \qquad
    y'\in Z_{r'-1-n+b+e(f)}^{s+n-b,t+n-b}(Y),
  \]
  for which
  \[
    d_{n-a'-b}^{\widehat f_{r'-1}}(\lambda^{a'}x')
       =\lambda^{n-b-e(f)}y'.
  \]
  This is a shorter essential contributor, and its target class is nonzero in
  the complete target coset measuring the discrepancy.  Consequently, if no
  such tuple exists, $\rho_*$ is surjective.  This is a statement about lifting
  representative solutions, not merely naturality of their page-level
  differentials.
\end{lemma}

\begin{proof}
  Represent an earlier solution by a filtered source lift together with its
  equality in the complete target coset.  Restrict the two-term filtered
  complexes for $\widehat f_{r-1}$ and $\widehat f_{r'-1}$.  Since the vertical
  $\rho$-ESS has only $d_0$, failure to lift this representative equation is
  exactly a class in the displayed cokernel; this defines
  $\operatorname{ob}_{r,r'}$.
  If that class is zero, alter the chosen equation by a restricted shorter
  correction and lift both sides, producing an element of
  $\operatorname{Sol}_{r}^{n}(x,y)$.  Restriction proves the reverse
  implication, and hence pointed exactness.

  For a nonzero obstruction, expand its target correction as a finite sum of
  shorter ESS images and choose a summand of least extension length that does
  not lift.  Replace an inessential summand by the higher-filtration essential
  contributor supplied by
  Proposition~\ref{prop:fextension-inessential-iff-higher}.  Length strictly
  decreases, so this process terminates.  The terminal source lifts at page
  $r'$ but not at page $r$, giving the displayed cycle-set difference.
  Filtration and weight preservation give the bounds on $a',b$ and the type of
  $y'$; minimality makes its target nonzero in the complete discrepancy coset.
\end{proof}

% SOURCE: aimpaper/main.tex:1981-2068, Proposition prop:dec738d3, corrected second part.
\begin{proposition}[Certificate for loss of essentiality on an earlier page]
  \label{prop:page-extension-loss-certificate}
  \uses{prop:page-extension-restrict,def:page-extension-crossing,def:page-extension-coset}
  \notready
  Suppose an essential $(f,E_r)$ extension of length $n$ becomes inessential
  on $E_{r'}$, where $2\le r'<r$ and
  $e(f)\le n\le r'-2+e(f)$.  Then there exist
  $0<a'\le n-e(f)$ and
  $0\le b\le n-a'-e(f)$, a source
  \[
    x'\in Z_{r'-1-a'}^{s+a',t+a'}(X)
      \setminus Z_{r-1-a'}^{s+a',t+a'}(X),
  \]
  and a typed target
  $y'\in Z_{r'-1-n+b+e(f)}^{s+n-b,t+n-b}(Y)$, nonzero in the complete target
  coset, such that
  \[
    d_{n-a'-b}^{\widehat f_{r'-1}}(\lambda^{a'}x')
       =\lambda^{n-b-e(f)}y'.
  \]
  Its target coset is one of the shorter images making the restricted target
  inessential.  This formulation supplies the ranges and target type omitted
  in the printed Proposition 6.20.
\end{proposition}

\begin{proof}
  Expand inessentiality in the earlier $\widehat f$-ESS and choose a shortest
  essential contributor.  If its source does not lift under $\rho$, it has the
  displayed cycle-set difference.  Otherwise lift it and replace it by the
  strictly shorter essential contributor responsible for the discrepancy.
  Finite descent reaches a nonliftable source.  Degree preservation fixes the
  target group and bounds $b$; essentiality puts $y'$ outside the complete
  target indeterminacy.
\end{proof}

% SOURCE: aimpaper/main.tex:1981-2081, the first-obstruction implication used to pass from Proposition prop:dec738d3 to Corollary cor:dfc6043e.
\begin{lemma}[First obstruction to lifting a page extension]
  \label{lem:page-extension-first-obstruction}
  \uses{def:page-extension,def:page-extension-coset,def:synthetic-ess-solution-tower,lem:page-extension-restriction-exact-obstruction}
  \notready
  Let $2\le r'<r<\infty$ and
  $e(f)\le n\le r'-2+e(f)$.  Suppose

  \[
    x\in Z_{r-1}^{s,t}(X),\qquad
    y\in Z_{r-1-n+e(f)}^{s+n,t+n}(Y),\qquad
    d_n^{f,E_{r'}}(x)=y.
  \]
  If the same relation $d_n^{f,E_r}(x)=y$ does not hold, then there are
  $0<a'\le n-e(f)$, $0\le b\le n-a'-e(f)$, and
  \[
    x'\in Z_{r'-1-a'}^{s+a',t+a'}(X)
       \setminus Z_{r-1-a'}^{s+a',t+a'}(X),
    \qquad
    y'\in Z_{r'-1-n+b+e(f)}^{s+n-b,t+n-b}(Y),
  \]
  such that
  \[
    d_{n-a'-b}^{\widehat f_{r'-1}}(\lambda^{a'}x')
       =\lambda^{n-b-e(f)}y'.
  \]
  The shorter contributor is essential and its target is nonzero in the
  complete target coset.  Thus failure of a finite lift produces the full
  obstruction tuple directly from the representative-restriction sequence.
\end{lemma}

\begin{proof}
  The $E_{r'}$ relation makes
  $\operatorname{Sol}_{r'}^n(x,y)$ nonempty; choose a solution $u$.  If $u$
  lay in the image of
  $\operatorname{Sol}_{r}^n(x,y)\to
  \operatorname{Sol}_{r'}^n(x,y)$, its lift would witness the asserted
  $E_r$ relation.  Hence its exact obstruction is nonzero.  Apply
  Lemma~\ref{lem:page-extension-restriction-exact-obstruction}; its nonzero
  obstruction representative is precisely the displayed tuple.
\end{proof}

% SOURCE: aimpaper/main.tex:2074-2081, Corollary cor:dfc6043e; the explicit
% inverse-limit obstruction hypothesis is a project correction for r=infinity.
\begin{corollary}[Page-extension stretching]
  \label{cor:page-extension-stretch}
  \uses{prop:page-extension-loss-certificate,lem:page-extension-first-obstruction,def:synthetic-ess-solution-tower,lem:synthetic-ess-coherent-tower-limit}
  \notready
  Let $r'\le r\le\infty$, suppose
  $e(f)\le n\le r'-2+e(f)$ and $d_n^{f,E_{r'}}(x)=y$, and suppose $x$
  and $y$ have the later-page types required by
  Definition~\ref{def:page-extension}: for finite $r$,
  \[
    x\in Z_{r-1}^{s,t}(X),\qquad
    y\in Z_{r-1-n+e(f)}^{s+n,t+n}(Y),
  \]
  while for $r=\infty$ both are permanent representatives.  If no tuple satisfying the loss-certificate conditions of
  Proposition~\ref{prop:page-extension-loss-certificate} exists between
  $r'$ and any finite page $q\le r$, then $d_n^{f,E_r}(x)=y$.
  When $r=\infty$, additionally assume that the finite representative-solution
  tower has vanishing torsor obstruction
  $\omega(\operatorname{Sol}_\bullet)$.  It is enough to supply a coherent
  solution tower, to prove every restriction on solution fibers surjective, or
  to verify the Mittag--Leffler condition of
  Lemma~\ref{lem:synthetic-ess-coherent-tower-limit}.
\end{corollary}

\begin{proof}
  For finite $r$, this is the contrapositive of the first-obstruction lemma.
  For $r=\infty$, the finite result makes every finite solution fiber
  nonempty.  The additional hypothesis and
  Lemma~\ref{lem:synthetic-ess-coherent-tower-limit} replace these unrelated
  choices by a coherent representative tower and identify its limit with a
  solution for $\widehat f$.  That solution is exactly
  $d_n^{f,E_\infty}(x)=y$ in the complete target coset.
\end{proof}

% SOURCE: aimpaper/main.tex:2084-2099, Example exam:stretchext; the solution-
% tower assertion makes explicit the representative compatibility used there.
\begin{proposition}[Stretching the Hopf-map extension to $E_\infty$]
  \label{prop:page-stretch-regression}
  \uses{prop:hopf-e3-page-extension,evidence:hopf-crossing-obstruction-exclusion,cor:page-extension-stretch,def:synthetic-ess-solution-tower,lem:page-extension-restriction-exact-obstruction,lem:synthetic-ess-coherent-tower-limit}
  \notready
  The $E_3$ extension
  $d_2^{\nu_{\mathrm{Hopf}},E_3}(h_0h_4^2)=h_0p$ has no loss certificate.
  Its finite representative-solution tower has surjective restriction maps,
  and therefore vanishing torsor obstruction.  Consequently it stretches to
  the essential extension
  $d_2^{\nu_{\mathrm{Hopf}},E_\infty}(h_0h_4^2)=h_0p$.
\end{proposition}

\begin{proof}
  The independent finite Hopf obstruction certificate exhausts the only
  degree-allowed shorter extension and excludes it by the Ext product relation
  at every finite quotient; its statement does not presuppose an
  $E_\infty$ relation.  It also supplies independently that
  $h_0h_4^2$ and $h_0p$ are permanent representatives, so both sides have the
  later-page types required by the stretching corollary.  Apply the finite
  exact-obstruction lemma to each adjacent pair:
  every representative solution lifts, so the solution-fiber restriction maps
  are surjective and the inverse-limit obstruction vanishes.  The source is a
  permanent cycle in the required range.  Apply page-extension stretching and
  the coherent-tower limit criterion.  Finally, the certificate's nonzero
  complete target-coset clause shows that the resulting length-two relation is
  essential.
\end{proof}

% SOURCE: aimpaper/main.tex:1398-1405 and 929-961; project packaging of the
% zero-homology rotation needed by the normalized-cofiber interface.
\begin{lemma}[Normalized Hopf cofiber exactness case]
  \label{lem:normalized-hopf-exactness-case}
  \uses{evidence:normalized-hopf-detection,def:normalized-synthetic-exactness-case,prop:positive-adams-filtration-homology-zero,def:cofiber-triangle,def:hf2-homology,thm:external-synthetic-triangle-lift}
  \notready
  The classical Hopf cofiber triangle
  \[
    S^3\xrightarrow{\nu_{\mathrm{Hopf}}}S^0\xrightarrow g
    S^0/\nu_{\mathrm{Hopf}}\xrightarrow hS^4
  \]
  carries the zero constructor of
  Definition~\ref{def:normalized-synthetic-exactness-case}.  Explicitly,
  \[
    H_*\nu_{\mathrm{Hopf}}=0,
    \qquad
    (e(\nu_{\mathrm{Hopf}}),e(g),e(h))=(1,0,0),
  \]
  and its long exact homology sequence supplies the rotated short exact
  sequence.  The constructor's lifted-triangle field is exactly the witness
  retained by Proposition~\ref{evidence:normalized-hopf-detection}, and its
  normalized Hopf component is therefore equal to $[h_2]$.
\end{lemma}

\begin{proof}
  The filtration-one detection by $h_2$ gives
  $e(\nu_{\mathrm{Hopf}})=1$, while degree considerations make the Hopf map
  zero on $H\mathbb F_2$-homology.  In the cofiber long exact sequence, $g$
  is the nonzero bottom-cell inclusion on homology and $h$ is the nonzero
  top-cell quotient.  The contrapositive of
  Proposition~\ref{prop:positive-adams-filtration-homology-zero} therefore
  gives Adams filtration zero and $e(g)=e(h)=0$.  The remaining
  long-exact-sequence terms split by degree into
  the required short exact sequence.  Package these facts and the displayed
  $e$-pattern around the precise lifted-triangle witness retained by the
  normalized-detection evidence; do not invoke the nonunique lift theorem a
  second time.  The evidence's component equality becomes the zero tag's
  equality $\widehat{\nu_{\mathrm{Hopf}}}=[h_2]$.
\end{proof}
